\documentclass[11pt,a4paper]{article}
\input{../../sablony/pripravy/priprava.sty}
\usepackage{tikz,pgfplots}
\pgfplotsset{compat=1.18}
\nastavPripravu{1. ročník}{Kvadratické rovnice a nerovnice}{112}{Soustava lineární a kvadratické rovnice}
\begin{document}
\blok{Soustava lineární a kvadratické rovnice}
Průsečíky přímky a paraboly, dosazovací postup a počet společných bodů.
\blok{Metoda a návaznost na graf}
\begin{itemize}
\item Vyjádři z lineární rovnice jednu proměnnou a dosazením získáš kvadratickou rovnici.
\item Dosaď všechny nalezené kořeny zpět a zapisuj uspořádané dvojice $[x;y]$.
\end{itemize}
\blok{Společně řešené příklady}
\begin{prikladReseniBox}[title={Příklad 1}]
$\begin{cases}y=x^{2}\\y=4\end{cases}$
\tcblower
\[
\begin{aligned}
x^{2}&=4\\[2pt]
0&=x^{2} - 4\\[2pt]
K&=\{[-2;4];\ [2;4]\}
\end{aligned}
\]
\end{prikladReseniBox}
\begin{prikladReseniBox}[title={Příklad 2}]
$\begin{cases}y=x^{2}\\y=2x + 3\end{cases}$
\tcblower
\[
\begin{aligned}
x^{2}&=2x + 3\\[2pt]
0&=x^{2} - 2x - 3\\[2pt]
K&=\{[-1;1];\ [3;9]\}
\end{aligned}
\]
\end{prikladReseniBox}
\begin{prikladReseniBox}[title={Příklad 3}]
$\begin{cases}y=x^{2} + 1\\y=2x\end{cases}$
\tcblower
\[
\begin{aligned}
x^{2} + 1&=2x\\[2pt]
0&=x^{2} - 2x + 1\\[2pt]
K&=\{[1;2]\}
\end{aligned}
\]
\end{prikladReseniBox}
\blok{Samostatné procvičování}
\begin{zadaniReseniBox}[title={Úloha 1}]
$\begin{cases}y=x^{2} - 2x\\y=0\end{cases}$
\tcblower
\[
\begin{aligned}
x^{2} - 2x&=0\\[2pt]
0&=x^{2} - 2x\\[2pt]
K&=\{[0;0];\ [2;0]\}
\end{aligned}
\]
\end{zadaniReseniBox}
\begin{zadaniReseniBox}[title={Úloha 2}]
$\begin{cases}y=4 - x^{2}\\y=x + 2\end{cases}$
\tcblower
\[
\begin{aligned}
4 - x^{2}&=x + 2\\[2pt]
0&=- x^{2} - x + 2\\[2pt]
K&=\{[-2;0];\ [1;3]\}
\end{aligned}
\]
\end{zadaniReseniBox}
\blok{Volitelná rezerva}
\poznamka{Studenti si vymyslí obdobný příklad se změnou jednoho koeficientu a vysvětlí, jak se změní počet kořenů či řešení.}
\blok{Volitelná výzva -- průsečíky přímky a paraboly}
\begin{zadaniReseniBox}[title={Rozšiřující příklad}]
$y=x^2-4$ a $y=2x-1$.
\tcblower
\[\begin{aligned}
x^2-4&=2x-1\\x^2-2x-3&=0\\x&=-1\ \text{nebo}\ x=3\\K&=\{[-1;-3];[3;5]\}
\end{aligned}\]
\end{zadaniReseniBox}
\blok{Domácí úkol}
\begin{zadaniReseniBox}[title={Domácí příprava}]
\begin{enumerate}[label=\alph*),itemsep=2pt]
\item $\begin{cases}y=x^{2} - 1\\y=3\end{cases}$
\item $\begin{cases}y=x^{2} - x\\y=x + 2\end{cases}$
\item $\begin{cases}y=x^{2}\\y=-1\end{cases}$
\item $\begin{cases}y=- x^{2} + 2x + 3\\y=2\end{cases}$
\item $\begin{cases}y=x^{2} - 4x + 3\\y=0\end{cases}$
\item $\begin{cases}y=2x^{2} - 3\\y=5\end{cases}$
\end{enumerate}
\tcblower
\begin{enumerate}[label=\alph*),itemsep=2pt]
\item
\[\begin{aligned}
x^{2} - 1&=3\\[2pt]
0&=x^{2} - 4\\[2pt]
K&=\{[-2;3];\ [2;3]\}
\end{aligned}\]
\item
\[\begin{aligned}
x^{2} - x&=x + 2\\[2pt]
0&=x^{2} - 2x - 2\\[2pt]
K&=\{[1 - \sqrt{3};3 - \sqrt{3}];\ [1 + \sqrt{3};\sqrt{3} + 3]\}
\end{aligned}\]
\item
\[\begin{aligned}
x^{2}&=-1\\[2pt]
0&=x^{2} + 1\\[2pt]
K&=\varnothing
\end{aligned}\]
\item
\[\begin{aligned}
- x^{2} + 2x + 3&=2\\[2pt]
0&=- x^{2} + 2x + 1\\[2pt]
K&=\{[1 - \sqrt{2};2];\ [1 + \sqrt{2};2]\}
\end{aligned}\]
\item
\[\begin{aligned}
x^{2} - 4x + 3&=0\\[2pt]
0&=x^{2} - 4x + 3\\[2pt]
K&=\{[1;0];\ [3;0]\}
\end{aligned}\]
\item
\[\begin{aligned}
2x^{2} - 3&=5\\[2pt]
0&=2x^{2} - 8\\[2pt]
K&=\{[-2;5];\ [2;5]\}
\end{aligned}\]
\end{enumerate}
\end{zadaniReseniBox}
\poznamka{Orientační čas: úvod 5 min, společné příklady 18 min, samostatně 15 min, kontrola a shrnutí 7 min. Rezerva je dobrovolná. Číselné parametry modelových aplikací nejsou skutečně naměřená data.}
\end{document}
